\documentclass[11pt]{article}
\usepackage[T1]{fontenc}
\usepackage{lmodern}
\usepackage[margin=1in,headheight=15pt]{geometry}
\usepackage{amsmath,amssymb,amsthm,mathtools}
\usepackage{microtype,booktabs,longtable,array,enumitem,xcolor,fancyhdr,xurl}
\usepackage[colorlinks=true,linkcolor=blue!40!black,citecolor=blue!40!black,urlcolor=blue!45!black]{hyperref}
\usepackage{bookmark}
\emergencystretch=2em
\setlength{\parskip}{3pt}
\numberwithin{equation}{section}
\newtheorem{theorem}{Theorem}[section]
\newtheorem{proposition}[theorem]{Proposition}
\newtheorem{lemma}[theorem]{Lemma}
\newtheorem{corollary}[theorem]{Corollary}
\newtheorem{conjecture}[theorem]{Conjecture}
\theoremstyle{definition}
\newtheorem{definition}[theorem]{Definition}
\theoremstyle{remark}
\newtheorem{remark}[theorem]{Remark}
\newtheorem{example}[theorem]{Example}
\DeclareMathOperator{\Li}{Li}
\DeclareMathOperator{\Cov}{Cov}
\DeclareMathOperator{\supp}{supp}
\DeclareMathOperator{\Impart}{Im}
\DeclareMathOperator{\Repart}{Re}
\DeclareMathOperator{\Log}{Log}
\newcommand{\dd}{\,\mathrm d}
\newcommand{\ii}{\mathrm i}
\newcommand{\R}{\mathbb R}
\newcommand{\C}{\mathbb C}
\newcommand{\N}{\mathbb N}
\newcommand{\E}{\mathbb E}
\newcommand{\cut}{\Omega}
\newcommand{\sv}{\boldsymbol s}
\newcommand{\xv}{\boldsymbol x}
\newcommand{\tv}{\boldsymbol t}
\newcommand{\pv}{\boldsymbol p}
\newcommand{\basecommit}{\nolinkurl{a0a90ef31877f98be437191c48b46f02d5456867}}
\pagestyle{fancy}
\fancyhf{}
\fancyhead[L]{\small Compensated positive kernels for polylogarithms}
\fancyhead[R]{\small ProveIt research}
\fancyfoot[C]{\thepage}
\hypersetup{pdftitle={Compensated Positive Kernels for Polylogarithms},pdfauthor={Research continuation prepared for ProveIt},pdfsubject={Subcritical angular zeros, sharp critical Euler constant, and exact generalized Stieltjes order}}
\title{\textbf{Compensated Positive Kernels\\for Polylogarithms}\\[8pt]
\large Subcritical zero geometry, a sharp critical Euler constant,\\and exact generalized Stieltjes order at every depth}
\author{Research continuation prepared for ProveIt\\\small With OpenAI assistance}
\date{9 October 2026 (America/Los\_Angeles)}
\begin{document}
\maketitle
\thispagestyle{empty}
\begin{abstract}
For $a,b>0$, put $F_{a,b}(z)=\sum_{n\ge1}H_{n-1}^{(b)}z^n/n^a$.
The preceding real-order report proves that its coefficients have a finite
signed Hausdorff representation precisely when $a+b\ge1$, and leaves
angular uniqueness below this boundary and a sharp critical Euler constant
as explicit conjectures. We prove both conjectures.

When $a+b\le1$, an explicit positive, generally infinite measure $\nu$
represents the coefficients through the compensated moments
$c_n=\int(1-u^{n-1})\dd\nu(u)$. Its finite first-difference measure gives
$F_{a,b}(z)=z^2(1-z)^{-1}\int(1-zu)^{-1}(1-u)\dd\nu(u)$.
This proves slit-plane nonvanishing below the old threshold, unique simple
angular zeros, and strict increase of the normalized radial displacement
throughout the closed subcritical region. Combined with the prior
supercritical theorem, nonvanishing and angular uniqueness hold for all
positive orders. A beta likelihood-ratio argument proves that
$-\Impart F_{a,w-a}(\ii\rho)$ decreases strictly with $a$ for every fixed
$w>0$ and $0<\rho\le1$. On $w=1$ this yields the conjectured optimal
uniform Euler constant $\pi/4+(\log2)/2$.

A separate positive transport formula supplies signed Euler differences
and explicit geometric remainders at every positive pair, including where
no finite signed measure exists. At arbitrary strict nesting depth $d$,
we construct a positive generalized Stieltjes representation of order $d$
and prove that this order is exact, using a logarithmic negative-axis
asymptotic. Equivalently, the coefficient normalization
$c_{n+d}n!/(\alpha)_n$ is a positive Hausdorff moment sequence exactly for
$\alpha\ge d$.

The package includes complete ordinary mathematical proofs, exact rational
certificates for nine Gaussian values and twelve subcritical angular
brackets, symbolic checks, and integration notes. The results are not
proof-assistant formalizations or independently refereed claims. The
$S_6$ period identity and the general supercritical normalized-radius
problem remain unresolved here.
\end{abstract}
\clearpage
\begingroup\small
\tableofcontents
\endgroup
\clearpage

\section{Scope, sources, and principal contributions}
\label{sec:scope}
\subsection{The concrete questions being continued}
This article continues the canonical manuscript \cite{repo} at revision
\basecommit, together with the incoming report
\emph{The Sharp Real-Order Threshold for Harmonic Polylogarithms}
\cite{threshold}. The latter is particularly important: its conjectures
are later than several already consolidated results. It explicitly asks
whether the angular zero remains unique when $a+b<1$, whether the
slit-plane nonvanishing theorem extends to that region, and whether
$a\mapsto-2\Impart F_{a,1-a}(\ii)$ decreases strictly with sharp supremum
$\pi/4+(\log2)/2$. These are the questions answered below.

The previous report already proves the finite signed-measure threshold,
the critical endpoint atom, critical normalized-radius monotonicity,
and the opposite behavior of $\Li_a(z)-z$. Its positive two-variable
kernel is also prior material. We use these insights, rederive the needed
formulas, and do not present them as discoveries of this article.
The canonical manuscript already has exact $S_4$ proofs and retains $S_6$
as a conjectural mixed period reduction. Neither status is changed here.

The inspection was targeted, not a line-by-line audit of every source
package. The incoming directory was inventoried through GitHub; the
threshold report was read from its accessible Library source. Its source
hash and the distinction between that copy and the repository archive
are recorded in \texttt{PROVENANCE.json}. The supplied package does not
modify a repository branch.

\begin{center}
\small
\begin{tabular}{@{}p{.25\linewidth}p{.31\linewidth}p{.35\linewidth}@{}}
\toprule
Question & Inspected prior result & Result proved here\\
\midrule
Finite signed moments & Exactly $a+b\ge1$ & Threshold retained; compensated positive moments fill $a+b<1$\\
Angular zero & Unique for $a+b\ge1$ & Unique and simple for all $a,b>0$\\
Principal-branch zeros & Only the double zero at zero for $a+b\ge1$ & Same statement for all $a,b>0$\\
Normalized radius & Increasing on $a+b=1$ & Increasing on all $a+b\le1$\\
Critical Euler constant & Explicit conjecture & Sharp constant $\pi/4+(\log2)/2$\\
All-order Euler evaluation & Finite-measure proof on $a+b\ge1$ & Positive transport proof at all positive orders\\
Strict depth $d$ & No such result in the inspected continuation & Exact Stieltjes order $d$; sharp normalization threshold\\
\bottomrule
\end{tabular}
\end{center}

\subsection{Conventions and the meaning of exactness}
Throughout,
\begin{equation}\label{eq:definitions}
 H_m^{(b)}=\sum_{j=1}^m j^{-b},\qquad
 c_n(a,b)=\frac{H_{n-1}^{(b)}}{n^a},\qquad
 F_{a,b}(z)=\sum_{n\ge1}c_n(a,b)z^n.
\end{equation}
We take $H_0^{(b)}=0$, so $c_1=0$ and $c_2=2^{-a}$.
The initial series is used on $|z|<1$ and continued to
\[\cut=\C\setminus[1,\infty).\]
All indices in this article are real and positive unless a boundary
value is expressly specified. Nested sums use strictly decreasing
indices, not the star convention. Depth-one branch conventions agree
with DLMF \S25.12 \cite{dlmf-polylog}. We abbreviate
\[
 w=a+b,\qquad g_{a,b}(\rho)=\Impart F_{a,b}(\ii\rho),\qquad
 g_{a,b}=g_{a,b}(1).
\]
Root-of-unity values on $w\le1$ mean analytic, equivalently Abel,
values. The original series then fails the ordinary term test.

A proved integral identity, a finite exact-arithmetic test, and a
floating-point diagnostic are different forms of evidence. Every theorem
below has a written proof. The accompanying tests check finite algebraic
identities, value enclosures, and implementation behavior; they do not
replace the proofs. No numerical agreement is used to prove a period
identity, and no statement of numerical period independence is made.
The novelty assessment is relative to the inspected project sources,
not a claim of exhaustive priority over the literature.

\section{Positive kernels and two different moment thresholds}
\label{sec:kernels}
Set
\begin{equation}\label{eq:k-h}
 k_s(T)=\frac{T^{s-1}}{\Gamma(s)},\qquad
 h(T)=\frac1{1-e^{-T}}-\frac1T\quad(T>0).
\end{equation}
The gamma integral gives
$\int_0^\infty e^{-qT}k_s(T)\dd T=q^{-s}$ for $q,s>0$.
The function $h$ increases strictly from $1/2$ to $1$, because
\[
 h'(T)=\frac1{T^2}-\frac1{4\sinh^2(T/2)}>0.
\]
Define
\begin{align}
 C_{a,b}(T)
 &=\int_0^T k_a(T-t)k_b(t)h(t)\dd t\label{eq:C-def}\\
 &=k_w(T)\E[h(TV)],\qquad V\sim\operatorname{Beta}(b,a).
 \label{eq:C-beta}
\end{align}
In particular $\tfrac12 k_w<C_{a,b}<k_w$.

\begin{lemma}[Regularized Hurwitz transform]\label{lem:hurwitz}
For $q>0$, $b>0$, $b\ne1$,
\begin{equation}\label{eq:hurwitz}
 \zeta(b,q)=\frac{q^{1-b}}{b-1}
       +\int_0^\infty e^{-qt}k_b(t)h(t)\dd t.
\end{equation}
Consequently
\begin{equation}\label{eq:C-laplace}
 \int_0^\infty e^{-qT}C_{a,b}(T)\dd T
 =q^{-a}\left(\zeta(b,q)-\frac{q^{1-b}}{b-1}\right).
\end{equation}
At $b=1$, the parenthesis is interpreted as $\log q-\psi(q)$.
\end{lemma}
\begin{proof}
For $b>1$, split the classical Hurwitz integral into $1/t+h(t)$.
The first part is $q^{1-b}/(b-1)$. Since $h$ is bounded, the other
integral is locally holomorphic for $\Repart b>0$. Analytic continuation
in $b$ proves the asserted formula away from its pole. This is the
regularization recorded in DLMF \S25.11 \cite{dlmf-hurwitz}.
Taking constant terms at $b=1$ gives $\log q-\psi(q)$.
Tonelli's theorem and the gamma integral prove the convolution formula.
\end{proof}

\begin{lemma}[A positive double kernel; prior construction]\label{lem:double}
For every $a,b>0$ and $z\in\cut$,
\begin{equation}\label{eq:double}
 F_{a,b}(z)=z^2\int_{0<y<x<1}
 \frac{W_{a,b}(x,y)}{(1-zx)(1-zy)}\dd y\dd x,
\end{equation}
where
\begin{equation}\label{eq:doubleweight}
 W_{a,b}(x,y)=\frac{(-\log x)^{a-1}(\log(x/y))^{b-1}}
 {\Gamma(a)\Gamma(b)}.
\end{equation}
This weight is positive and has total mass $2^{-a}$.
\end{lemma}
\begin{proof}
For nonnegative integers $j,k$, substitution $y=xv$ gives
\[
 \int W_{a,b}(x,y)x^jy^k\dd y\dd x
 =\frac1{(j+k+2)^a(k+1)^b}.
\]
Expanding the two denominators in the disk and summing over $j+k=n-2$
therefore gives $c_n$. Absolute convergence justifies the expansion.
On a compact subset of $\cut$, $1-zx$ and $1-zy$ stay uniformly away
from zero for $0\le x,y\le1$. The finite integral is holomorphic there
and supplies the continuation. Take $j=k=0$ for its mass.
\end{proof}

The following preceding result is included to make the all-order
conclusions independent of an unquoted sign argument.
\begin{proposition}[The finite signed threshold; prior theorem]
\label{prop:signedthreshold}
There is a finite real signed measure $\mu$ on $[0,1]$ satisfying
$c_n=\int u^{n-1}\dd\mu(u)$ for all $n\ge1$ if and only if $w\ge1$.
For $w>1$ it has density $K_{a,b}(-\log u)$, where
\begin{align}
 K_{a,b}(T)&=\zeta(b)k_a(T)-\frac{k_{w-1}(T)}{b-1}-C_{a,b}(T)
 &&(b\ne1),\label{eq:superkernel}\\
 K_{a,1}(T)&=k_a(T)(\gamma+\psi(a)-\log T)-C_{a,1}(T).
 \label{eq:superkernel1}
\end{align}
This density has exactly one sign change, negative to positive as $u$
increases. On $w=1$, the measure instead has the form
$\mu=a^{-1}\delta_1-\nu$ with $\nu$ a positive finite measure.
\end{proposition}
\begin{proof}
The densities in \eqref{eq:superkernel}--\eqref{eq:superkernel1} are
integrable with weight $e^{-T}\dd T$. Their Laplace transforms are
$c_n$ by Lemma~\ref{lem:hurwitz}; at $b=1$, also use
\[
 \int_0^\infty e^{-qT}k_a(T)\log T\dd T
 =q^{-a}(\psi(a)-\log q).
\]
In particular the signed measure has total mass zero.
For $b\ne1$, differentiation after dividing by $k_a(T)$ gives
\begin{align*}
 \left(\frac{K_{a,b}}{k_a}\right)'(T)
 ={}&-\frac{\Gamma(a)}{\Gamma(w-1)}T^{b-2}\\
 &-\frac{\Gamma(a)}{\Gamma(w)}T^{b-1}
       \E[b h(TV)+TVh'(TV)]<0.
\end{align*}
For $b=1$ the derivative is
$-T^{-1}-a^{-1}\E[h(TV)+TVh'(TV)]<0$.
The normalized density starts positive (possibly at $+\infty$) and
tends to $-\infty$. This proves exactly one transversal crossing, with
the stated orientation after $u=e^{-T}$.

The construction on $w=1$ follows from
Theorem~\ref{thm:compensated} below, where $\nu$ has mass $1/a$.
For $w<1$, the integral test gives
$c_n\sim n^{1-w}/(1-b)\to\infty$, impossible for moments of a finite
signed measure on $[0,1]$. Uniqueness in the finite cases follows from
density of polynomials in the continuous functions on $[0,1]$.
\end{proof}

\section{The compensated subcritical measure}
\label{sec:compensated}
The absence of a finite signed measure does not preclude a positive
representation of \emph{differences}. Here that distinction is decisive.

\begin{theorem}[Explicit positive compensation]\label{thm:compensated}
Let $a,b>0$ and $w=a+b\le1$. Put
\[
 \delta=1-w\in[0,1),\qquad A=\frac1{1-b},\qquad
 N_{a,b}(T)=-\zeta(b)k_a(T)+C_{a,b}(T).
\]
Define a positive measure on $(0,1)$ by
\begin{equation}\label{eq:nu}
 \dd\nu_{a,b}(u)=
 \left[N_{a,b}(-\log u)
 +\frac{\delta}{(1-b)\Gamma(1-\delta)}
             (-\log u)^{-1-\delta}\right]\dd u.
\end{equation}
When $\delta=0$, the second summand is omitted. Then
\begin{equation}\label{eq:compensatedmoments}
 c_n=\int_0^1(1-u^{n-1})\dd\nu_{a,b}(u),\qquad n\ge1,
\end{equation}
and
\begin{equation}\label{eq:omega}
 \dd\omega_{a,b}(u)=(1-u)\dd\nu_{a,b}(u),\qquad
 \omega_{a,b}((0,1))=2^{-a}.
\end{equation}
The density of $\omega_{a,b}$ is strictly positive throughout $(0,1)$.
The measure $\nu_{a,b}$ is finite precisely on $w=1$, where its mass
is $A=1/a$. For $w<1$, its total mass is infinite, but every compensated
integral in \eqref{eq:compensatedmoments} is finite.
\end{theorem}
\begin{proof}
Since $0<b<1$, $\zeta(b)<0$. For example, the alternating Dirichlet
series is positive and $\zeta(b)=\eta_D(b)/(1-2^{1-b})$.
Thus $N_{a,b}>0$. Lemma~\ref{lem:hurwitz} gives, for $q>0$,
\begin{equation}\label{eq:qcomp}
 q^{-a}(\zeta(b)-\zeta(b,q))
   =Aq^\delta-\int_0^\infty e^{-qT}N_{a,b}(T)\dd T.
\end{equation}
At $q=1$ the left side is zero, so
$\int e^{-T}N_{a,b}(T)\dd T=A$.
For $0<\delta<1$, integration by parts gives
\begin{equation}\label{eq:fractionalpower}
 q^\delta-1=\frac{\delta}{\Gamma(1-\delta)}
 \int_0^\infty(e^{-T}-e^{-qT})T^{-1-\delta}\dd T.
\end{equation}
Indeed the boundary term vanishes at both endpoints, and the remaining
gamma integrals evaluate the expression. Subtract the $q=1$ identity
from \eqref{eq:qcomp}, insert \eqref{eq:fractionalpower}, and change
variables $u=e^{-T}$. For integer $q=n$, the Hurwitz difference is
$H_{n-1}^{(b)}$, proving \eqref{eq:compensatedmoments}. At $\delta=0$
the same argument works without the fractional-power term.

Near $u=1$, the new density has leading behavior
\[
 \frac{\delta}{(1-b)\Gamma(1-\delta)}(1-u)^{-1-\delta}
 \qquad(\delta>0).
\]
Its uncompensated mass diverges, while multiplication by $1-u$ makes
it integrable because $\delta<1$. The remaining density is integrable:
its $e^{-T}$-weighted mass is $A$. At $u=0$ exponential decay in the
$T$ coordinate makes all terms integrable. Setting $n=2$ gives the
mass in \eqref{eq:omega}. Strict positivity is immediate.
\end{proof}

\begin{remark}[The cancellation is mathematical, not notational]
For $w<1$, it is invalid to rewrite
$\int(1-u^{n-1})\dd\nu$ as the difference of its two separate integrals:
both are infinite. Nor is $\infty\delta_1-\nu$ a finite signed measure.
The finite object to use for zero geometry is $\omega=(1-u)\nu$.
This retains, rather than contradicts, the threshold of
Proposition~\ref{prop:signedthreshold}.
\end{remark}

\begin{corollary}[A positive first-difference Stieltjes factor]
\label{cor:factor}
For $w\le1$ and $z\in\cut$,
\begin{equation}\label{eq:factor}
 F_{a,b}(z)=\frac{z^2}{1-z}P_{a,b}(z),\qquad
 P_{a,b}(z)=\int_0^1\frac{\dd\omega_{a,b}(u)}{1-zu}.
\end{equation}
Equivalently,
\begin{equation}\label{eq:differencegf}
 \frac{(1-z)F_{a,b}(z)}{z^2}
 =\sum_{n\ge0}(c_{n+2}-c_{n+1})z^n.
\end{equation}
The removable value at zero is $2^{-a}$.
\end{corollary}
\begin{proof}
In the disk,
\[
 \sum_{n\ge1}(1-u^{n-1})z^n
 =\frac{z^2(1-u)}{(1-z)(1-zu)}.
\]
Termwise integration is justified by
$|1-u^{n-1}|\le(n-1)(1-u)$ and finiteness of $\omega$.
Both sides continue holomorphically to $\cut$, proving the claim.
\end{proof}

\begin{theorem}[A sharp complementary positivity threshold]
\label{thm:alternation}
For $a,b>0$, the sequence $(c_{n+2}-c_{n+1})_{n\ge0}$ is a positive
Hausdorff moment sequence if and only if $a+b\le1$. In that region,
with $\nabla f_n=f_{n+1}-f_n$,
\begin{equation}\label{eq:alternation}
 (-1)^{k-1}\nabla^k c_{n+1}
 =\int_0^1u^n(1-u)^k\dd\nu_{a,b}(u)>0
 \quad(n\ge0,\ k\ge1).
\end{equation}
All finite principal Hankel matrices of the first-difference sequence
are positive definite. The continuous function
\begin{equation}\label{eq:Bernstein}
 B_{a,b}(x)=(x+1)^{-a}\bigl(\zeta(b)-\zeta(b,x+1)\bigr),\qquad x\ge0,
\end{equation}
is a Bernstein function: it is nonnegative and its derivative is
strictly completely monotone.
\end{theorem}
\begin{proof}
Equation~\eqref{eq:alternation} follows by taking finite differences
inside \eqref{eq:compensatedmoments}. Its integrals are finite because
$k\ge1$. The case $k=1$ identifies the first-difference sequence as
the moments of $\omega$. For a nonzero polynomial $p$,
$\int p(u)^2\dd\omega(u)>0$, since $\omega$ has positive density on
every interior interval. This proves strict Hankel positivity.

When $w>1$, $c_n\to0$: for $b<1$ use the integral test, for $b=1$
use $H_{n-1}=O(\log n)$, and for $b>1$ use boundedness of the harmonic
sums. Since $c_1=0<c_2$, the differences cannot all be nonnegative.
They therefore cannot be moments of a positive measure. This proves
necessity, not just a sufficient construction.

The proof of Theorem~\ref{thm:compensated} also works for $q=x+1$,
giving $B(x)=\int(1-u^x)\dd\nu(u)$. Differentiation for $x>0$ yields
\[
 (-1)^j B^{(j+1)}(x)
 =\int_0^1u^x(-\log u)^{j+1}\dd\nu(u)>0.
\]
Near one the integrability condition is again $\delta<1$, and near
zero the $u$ density supplies exponential decay. The same domination
also gives continuity at $x=0$.
\end{proof}

\section{All-positive-order nonvanishing and angular uniqueness}
\label{sec:zeros}
\begin{theorem}[No nonzero zeros on the principal slit plane]
\label{thm:zerofree}
For every $a,b>0$, the only zero of $F_{a,b}$ on $\cut$ is its double
zero at $0$.
\end{theorem}
\begin{proof}
For $w\le1$, use \eqref{eq:factor}. If $\Impart z\ne0$, then
\[
 \Impart P_{a,b}(z)
 =\Impart z\int_0^1\frac{u\dd\omega_{a,b}(u)}{|1-zu|^2},
\]
whose integral is strictly positive. For real $z<1$, $P_{a,b}(z)>0$.
The factor $1-z$ has no zero on $\cut$, and $P_{a,b}(0)=2^{-a}$.

For completeness, when $w>1$, let $r\in(0,1)$ be the crossing of
Proposition~\ref{prop:signedthreshold}. The positive nonzero measure
$\lambda=(u-r)\mu$ gives
\[
 F_{a,b}(z)=\frac{z^2}{1-rz}
           \int_0^1\frac{\dd\lambda(u)}{1-zu}.
\]
This follows by subtracting $z(1-rz)^{-1}\int\dd\mu=0$ from its
finite signed-measure representation. The same imaginary-part test
proves nonvanishing of the integral, and $1-rz$ has no zero on $\cut$.
The coefficient $c_2=2^{-a}$ gives exact multiplicity two.
\end{proof}

\subsection{A root lemma that permits infinite uncompensated mass}
\begin{lemma}[Endpoint factor and normalized root]\label{lem:root}
Let $\omega$ be a finite positive measure on $(0,1)$ with positive
density there. Define
\[
 F(z)=\frac{z^2}{1-z}\int_0^1\frac{\dd\omega(u)}{1-zu}.
\]
For every $0<\rho\le1$, $\Impart F(\rho e^{\ii\theta})$ has exactly
one zero in $0<\theta<\pi$. It is simple, has positive sign before
and negative sign after the zero, and satisfies
\begin{equation}\label{eq:stronganglerange}
 \arccos\rho<\theta(\rho)<\arccos(\rho/2).
\end{equation}
Moreover $\eta(\rho)=\cos\theta(\rho)/\rho$ increases strictly.
\end{lemma}
\begin{proof}
Write
\[
 x=\rho^2,\qquad c=\cos\theta,\qquad
 m=\frac{2c}{\rho}-1,\qquad
 D(u)=1+x\bigl(u^2-(1+m)u\bigr).
\]
For the relevant root range $0<m<1$, all denominators are positive,
including $D(1)=1-xm$. Direct rationalization gives
\begin{equation}\label{eq:angularG}
 \Impart F(\rho e^{\ii\theta})
 =\frac{\rho^3\sin\theta}{1-xm}G(m,x),\qquad
 G(m,x)=\int_0^1\frac{m-u}{D(u)}\dd\omega(u).
\end{equation}
At $m=0$, $G<0$. At $m=1$ and $x<1$, $G>0$. Its derivative is
\begin{equation}\label{eq:Gm}
 G_m(m,x)=\int_0^1\frac{1-xu}{D(u)^2}\dd\omega(u)>0.
\end{equation}
Therefore there is exactly one root $0<m<1$ at every $x<1$.

At $x=1$, an endpoint argument is essential. Put $\varepsilon=1-m$.
Then $D(u)=(1-u)^2+\varepsilon u$. On the negative part $u>m\ge1/2$,
\[
 \left|\frac{m-u}{D(u)}\right|\le\frac1u\le2.
\]
Its integral tends to zero, since $\omega((m,1))\to0$. On any fixed
interior interval below $m$, the integrand tends to $(1-u)^{-1}>0$.
Thus $G(m,1)>0$ for $m$ sufficiently close to one. This argument does
not require $\int(1-u)^{-1}\dd\omega$ to be finite; in the subcritical
application it is infinite. Equation~\eqref{eq:Gm} now proves the
unique root at $x=1$ as well.

For actual angular parameters, the numerator $2c-\rho(1+u)$ is
negative when $c\le\rho/2$ and positive when $c\ge\rho$, whenever
these ranges occur. Inside this interval $G$ increases with $m$,
whereas $m$ decreases with $\theta$. This proves the global sign
pattern, the strict range, and simplicity. Smooth dependence follows
from the implicit function theorem, since the root stays away from
$z=1$ even when $\rho=1$.

For the radial sign, set $d(u)=u^2-(1+m)u$. At a root of $G$, the
probability measure proportional to $D(u)^{-1}\dd\omega(u)$ has mean
$m$. Since
\[
 d(m)=-m,\qquad d(u)-d(m)=(u-m)(u-1),
\]
subtracting the constant $d(m)/(1-xm)$ in the expression for $G_x$
gives
\[
 G_x=-\frac1{1-xm}\int_0^1
              \frac{(u-m)^2(1-u)}{D(u)^2}\dd\omega(u)<0.
\]
It follows that
\begin{equation}\label{eq:radialderivative}
 \frac{\dd m}{\dd x}=
 \frac{\displaystyle\int_0^1
              \frac{(u-m)^2(1-u)}{D(u)^2}\dd\omega(u)}
 {(1-xm)\displaystyle\int_0^1
              \frac{1-xu}{D(u)^2}\dd\omega(u)}>0.
\end{equation}
Finally $\eta=(1+m)/2$, so $\eta'(\rho)=\rho m'(\rho^2)>0$.
Strict positivity follows from the nondegenerate interior density.
\end{proof}

\begin{theorem}[Resolution of subcritical angular uniqueness]
\label{thm:subcriticalangular}
For all $a,b>0$ with $a+b\le1$ and $0<\rho\le1$, the unique upper
semicircle zero $\theta_{a,b}(\rho)$ satisfies
\eqref{eq:stronganglerange}. Its normalized displacement
\[
 \eta_{a,b}(\rho)=\frac{\cos\theta_{a,b}(\rho)}\rho
\]
increases strictly and lies in $(1/2,1)$. Formula
\eqref{eq:radialderivative}, with $\omega=\omega_{a,b}$, is an exact
positive expression for the derivative.
\end{theorem}
\begin{proof}
Apply Lemma~\ref{lem:root} to the finite positive measure in
Corollary~\ref{cor:factor}.
\end{proof}

This proves the conjecture labeled \texttt{conj:subcritical} in
\cite{threshold}, including unit radius and simplicity. The extension
of normalized-radius monotonicity from the critical line to the entire
subcritical triangle is an additional conclusion.

\begin{corollary}[Angular uniqueness at every positive pair]
\label{cor:allangular}
For all $a,b>0$ and $0<\rho\le1$, there is exactly one simple zero of
$\Impart F_{a,b}(\rho e^{\ii\theta})$ in $0<\theta<\pi$. The sign
is positive before it and negative after it, and
\[
 \arccos\rho<\theta_{a,b}(\rho)<\pi/2.
\]
\end{corollary}
\begin{proof}
Only $w>1$ remains. Let $\mu$ be the one-crossing finite measure in
Proposition~\ref{prop:signedthreshold}, and set
\[
 K_{\rho,c}(u)=\frac1{1-2\rho cu+\rho^2u^2},\qquad
 J(\rho,c)=\int K_{\rho,c}\dd\mu.
\]
The imaginary part equals $\rho\sin\theta\,J(\rho,c)$.
An increasing test function integrates positively against a nonzero
zero-mass measure whose negative support lies to the left of its
positive support: subtract its value at the crossing. Thus
$J(\rho,0)<0$ and $J(\rho,\rho)>0$ for $\rho<1$.
At $\rho=1$, the negative part is supported away from one and
\[
 \lim_{c\uparrow1}J(1,c)
 =\int(1-u)^{-2}\dd\mu(u)=\sum_{n\ge1}n c_n>0,
\]
with $+\infty$ allowed. Separate the signs before applying monotone
convergence or expanding the nonnegative polynomial partial sums.

For $c'>c$ the ratio $K_{\rho,c'}/K_{\rho,c}$ increases strictly on
$(0,1)$, since its derivative is
\[
 \frac{2\rho(c'-c)(1-\rho^2u^2)}
 {(1-2\rho c'u+\rho^2u^2)^2}>0.
\]
At a zero of $J$, apply the same ordered-sign test to
$K_{\rho,c}\mu$, which again has zero mass. This proves uniqueness
and the global sign pattern. At the zero,
\[
 J_c=\int K_{\rho,c}(u)
          \frac{2\rho u}{1-2\rho cu+\rho^2u^2}\dd\mu(u)>0,
\]
because the second factor increases strictly. The angular derivative
is $-\rho\sin^2\theta J_c<0$. This is the prior supercritical argument,
now combined with the newly proved subcritical case.
\end{proof}

\begin{corollary}[Strict small-radius coefficient below the threshold]
\label{cor:smallradius}
For $w\le1$,
\begin{equation}\label{eq:smallradius}
 \eta_{a,b}(\rho)=\frac{c_3}{2c_2}
 +\left(\frac{c_3c_4}{c_2^2}-\frac{c_5}{2c_2}
                    -\frac{c_3^3}{2c_2^3}\right)\rho^2+O(\rho^4).
\end{equation}
The coefficient of $\rho^2$ is strictly positive.
\end{corollary}
\begin{proof}
At $x=0$, $m_0=(c_3-c_2)/c_2$ is the mean of $U$ with law
$\omega/c_2$. Formula~\eqref{eq:radialderivative} gives the coefficient
$\tfrac12\E[(U-m_0)^2(1-U)]>0$.
The moments $\int u^j\dd\omega=c_{j+2}-c_{j+1}$ expand it into
\eqref{eq:smallradius}. Analytic implicit dependence near $x=0$
justifies the remainder. The expansion itself is not new; the strict
sign on the whole subcritical region is.
\end{proof}

\section{A fixed-weight Gaussian monotonicity theorem}
\label{sec:gaussian}
\begin{theorem}[Strict monotonicity along every fixed total weight]
\label{thm:weightmonotonicity}
Fix $w>0$ and $0<\rho\le1$. The function
\begin{equation}\label{eq:fixedweight}
 a\longmapsto-\Impart F_{a,w-a}(\ii\rho)
\end{equation}
is strictly decreasing on $0<a<w$.
\end{theorem}
\begin{proof}
In \eqref{eq:double}, put $x=e^{-s}$, $y=e^{-(s+t)}$ and $T=s+t$.
For $X=\rho e^{-s}$ and $Y=\rho e^{-T}$, define
\begin{equation}\label{eq:Q}
 Q_\rho(s,T)=\frac{XY(X+Y)}{(1+X^2)(1+Y^2)}.
\end{equation}
Then
\begin{align}
 -\Impart F_{a,b}(\ii\rho)
 &=\frac1{\Gamma(a)\Gamma(b)}
   \int_0^\infty\int_0^\infty
          s^{a-1}t^{b-1}Q_\rho(s,s+t)\dd t\dd s\notag\\
 &=\frac1{\Gamma(w)}\int_0^\infty T^{w-1}
       \E[Q_\rho(TV,T)]\dd T,
 \quad V\sim\operatorname{Beta}(a,w-a).
 \label{eq:beta-gaussian}
\end{align}
For fixed $T$, the function $s\mapsto Q_\rho(s,T)$ decreases strictly.
Indeed $X$ decreases with $s$, and
\begin{equation}\label{eq:Qderivative}
 \frac{\partial}{\partial X}
 \frac{XY(X+Y)}{(1+X^2)(1+Y^2)}
 =\frac{Y\{2X+Y(1-X^2)\}}
 {(1+X^2)^2(1+Y^2)}>0,
\end{equation}
where $0<X\le\rho\le1$.

For $a_2>a_1$, the ratio of the two beta densities is a positive
constant times $(v/(1-v))^{a_2-a_1}$, which increases strictly. The
expectation of any strictly decreasing function therefore decreases
strictly. One direct proof differentiates in $a$:
\begin{equation}\label{eq:covariance}
 \frac{\partial}{\partial a}\E[q(V)]
 =\Cov\left(q(V),\log\frac V{1-V}\right)<0.
\end{equation}
The sign follows by writing covariance as half the expectation of
$(q(V)-q(V'))(L(V)-L(V'))$ for independent $V,V'$; the factors have
opposite signs off the diagonal. Apply this with $q(v)=Q_\rho(Tv,T)$.
Since $Q_\rho(Tv,T)\le2\rho^3e^{-T}$, the $T$ integral is dominated.
Beta logarithmic moments are locally uniformly finite for $0<a<w$,
so differentiation and integration can be interchanged as well.
Strictness survives integration over every $T>0$.
\end{proof}

\subsection{Explicit endpoint identities and inequalities}
Write $\beta(s)$ for the analytically continued Dirichlet beta function,
and $\eta_D(s)=(1-2^{1-s})\zeta(s)$ for the Dirichlet eta function,
with its removable value at $s=1$. This notation is distinct from the
normalized angular displacement $\eta_{a,b}(\rho)$.

\begin{corollary}[Sharp fixed-weight Gaussian bounds]\label{cor:weightbounds}
For $0<a<w$,
\begin{equation}\label{eq:weightbounds}
 \beta(w)-\beta(w-1)<-g_{a,w-a}
 <\frac{\beta(w)+2^{-w}\eta_D(w)}2.
\end{equation}
Both endpoints are sharp limits, attained as $a\uparrow w$ and
$a\downarrow0$, respectively, but neither is attained at an interior
parameter. For every $w>0$ and every $a\in(0,w)$,
\begin{equation}\label{eq:universalgaussian}
 -2g_{a,w-a}<\frac{1+\sqrt2}{2}.
\end{equation}
\end{corollary}
\begin{proof}
As $a\downarrow0$, the beta variable in \eqref{eq:beta-gaussian}
concentrates at zero; as $a\uparrow w$, it concentrates at one.
The preceding domination proves convergence of the expectations and
integrals. The limiting analytic functions are
\begin{equation}\label{eq:limitidentities}
 F_{0,w}(z)=\frac z{1-z}\Li_w(z),\qquad
 F_{w,0}(z)=\Li_{w-1}(z)-\Li_w(z).
\end{equation}
These identities follow first by coefficient summation in the disk,
and their integral continuations give the endpoint values at $z=\ii$.
At that point,
\[
 \Li_w(\ii)=-2^{-w}\eta_D(w)+\ii\beta(w).
\]
This yields \eqref{eq:weightbounds}. When $w\le1$, the occurrence of
$\beta(w-1)$ is an analytic value, not the sum of a convergent ordinary
Dirichlet series.

Twice the upper limit in \eqref{eq:weightbounds} also equals
\[
 \frac1{\Gamma(w)}\int_0^\infty T^{w-1}e^{-T}
                \frac{1+e^{-T}}{1+e^{-2T}}\dd T.
\]
The scalar function $(1+y)/(1+y^2)$ on $0<y<1$ has maximum
$(1+\sqrt2)/2$, at $y=\sqrt2-1$. Since the gamma measure is not
concentrated there, its expectation is strictly smaller.
\end{proof}

\begin{remark}[Why the restriction on the Gaussian radius matters]
The derivative in \eqref{eq:Qderivative} is uniformly positive in the
present argument because $X\le1$. For radii above one, $Y(1-X^2)$ can
have the other sign. The theorem does not assert monotonicity for all
$\rho>1$, and the proof cannot be extended there by dropping a hypothesis.
\end{remark}

\section{Euler acceleration without a finite signed measure}
\label{sec:euler}
\subsection{A bounded order-one transport functional}
Let
\begin{equation}\label{eq:transportmeasures}
 \dd\sigma_a(v)=k_a(-\log v)\dd v,\qquad
 \dd\tau_b(u)=\frac{k_b(-\log u)}{1-u}\dd u\quad(0<u,v<1).
\end{equation}
The measure $\sigma_a$ is a probability measure, while $\tau_b$ may
have infinite mass. Nevertheless
$\int(1-u)\dd\tau_b(u)=1$.

\begin{theorem}[All-order positive transport]\label{thm:transport}
For all $a,b>0$, the functional
\begin{equation}\label{eq:L}
 \mathcal L_{a,b}(\varphi)
 =\int_0^1\int_0^1\bigl(\varphi(v)-\varphi(uv)\bigr)
                \dd\sigma_a(v)\dd\tau_b(u)
\end{equation}
is well-defined for $\varphi\in C^1[0,1]$, and
\begin{equation}\label{eq:Lnorm}
 |\mathcal L_{a,b}(\varphi)|\le2^{-a}\|\varphi'\|_\infty.
\end{equation}
The constant is exact. Its polynomial moments are
$\mathcal L_{a,b}(u^{n-1})=c_n$, and
\begin{equation}\label{eq:transportF}
 F_{a,b}(z)=\mathcal L_{a,b}\left(u\longmapsto\frac z{1-zu}\right)
 \quad(z\in\cut).
\end{equation}
When $a+b<1$, this functional cannot satisfy a bound by
$C\|\varphi\|_\infty$: its distributional order is exactly one, not zero.
\end{theorem}
\begin{proof}
The mean-value estimate is
$|\varphi(v)-\varphi(uv)|\le v(1-u)\|\varphi'\|_\infty$.
Its integral is $2^{-a}\|\varphi'\|_\infty$; equality occurs for
$\varphi(u)=u$. For a monomial,
\[
 \mathcal L(u^{n-1})
 =\left(\int v^{n-1}\dd\sigma_a(v)\right)
   \left(\int\frac{1-u^{n-1}}{1-u}k_b(-\log u)\dd u\right)
 =n^{-a}H_{n-1}^{(b)}.
\]
The power series of $z/(1-zu)$ converges in $C^1[0,1]$ for $|z|<1$,
so the moment equality proves \eqref{eq:transportF} there.
On compact slit sets, the same test functions and their derivatives
are uniformly bounded; hence their images under $\mathcal L$ are
holomorphic and continuation proves the formula everywhere.
For $w<1$, the test polynomials $u^{n-1}$ have supremum norm one but
$\mathcal L(u^{n-1})=c_n\to\infty$, excluding any order-zero bound.
\end{proof}

\subsection{Strict Euler signs and explicit remainders}
For an integer stride $h\ge1$, put
\begin{equation}\label{eq:dk}
 d_k^{(h)}=\sum_{n=0}^k(-1)^n\binom kn c_{hn+1}.
\end{equation}
\begin{theorem}[All-positive-order Euler differences]\label{thm:Eulerall}
For every $a,b>0$ and $h\ge1$,
\begin{equation}\label{eq:differencebound}
 d_0^{(h)}=0,\qquad
 0<-d_k^{(h)}<H_h^{(b)}\le h\quad(k\ge1).
\end{equation}
More precisely,
\begin{equation}\label{eq:positivedifference}
 -d_k^{(h)}=\iint\left[(1-u^hv^h)^k-(1-v^h)^k\right]
                    \dd\sigma_a(v)\dd\tau_b(u).
\end{equation}
For the Gaussian stride $h=2$, let
\begin{equation}\label{eq:EN}
 E_N=\sum_{k=0}^{N-1}\frac{d_k^{(2)}}{2^{k+1}}\qquad(N\ge1).
\end{equation}
Then $E_N$ decreases strictly to the analytic value $g_{a,b}$, and
\begin{equation}\label{eq:allEulerbound}
 E_N-H_2^{(b)}2^{-N}<g_{a,b}<E_N.
\end{equation}
The exact positive tail identity is
\begin{equation}\label{eq:allEulertail}
 2^N(E_N-g_{a,b})
 =\iint\bigl[q_N(uv)-q_N(v)\bigr]
                    \dd\sigma_a(v)\dd\tau_b(u),
 \qquad q_N(v)=\frac{(1-v^2)^N}{1+v^2}.
\end{equation}
\end{theorem}
\begin{proof}
The binomial theorem inside the finite sum gives
\eqref{eq:positivedifference}. It is strictly positive.
For $U=u^h$, $V=v^h$, convexity of $f(t)=(1-t)^k$ gives
\[
 0<f(UV)-f(V)\le(1-U)\{1-f(V)\}<1-U.
\]
Integrate and use
$\int(1-u^h)\dd\tau_b(u)=H_h^{(b)}$. This also justifies every
integral despite a possibly infinite mass of $\tau_b$.

The Euler series is absolutely convergent by the uniform bound on
$d_k^{(2)}$. Tonelli applies to the positive integrands for $-d_k$.
Summing them gives
\[
 \sum_{k\ge0}\frac{d_k^{(2)}}{2^{k+1}}
 =\iint\left(\frac1{1+v^2}-\frac1{1+u^2v^2}\right)
                  \dd\sigma_a(v)\dd\tau_b(u)
 =g_{a,b},
\]
where the last equality is \eqref{eq:transportF} at $z=\ii$.
Summing the remaining geometric series proves
\eqref{eq:allEulertail}. Since $q_N$ decreases strictly, its integrand
is positive. Finally sum the strict coefficient bound over $k\ge N$
to get \eqref{eq:allEulerbound}.
\end{proof}

The estimate $H_2^{(b)}<2$ is a convenient rational certificate bound
at arbitrary positive orders. It is not claimed to be the optimal
all-order Gaussian error constant. The preceding constant one for
$a\ge1$ remains sharper on its stated domain and is not superseded by
a weaker estimate. The coefficient bounds themselves have sharp
uniform constants: for $h=1$, let $a\downarrow0$, $b\to\infty$ to
obtain $d_k^{(1)}\to-1$ for fixed $k\ge1$; for $h=2$, already
$-d_1^{(2)}=H_2^{(b)}/3^a\to2$ as $a,b\downarrow0$.

\begin{corollary}[Independent affine strides]\label{cor:affine}
Let $a,b>0$, let $r$ be a positive integer, and let $\lambda>0$.
For $f(n)=H_{rn}^{(b)}/(\lambda n+1)^a$, define
$d_k=\sum_{n=0}^k(-1)^n\binom kn f(n)$. Then $d_0=0$ and
$0<-d_k<H_r^{(b)}$ for every $k\ge1$. The Euler sums
$E_N(f)=\sum_{k<N}d_k/2^{k+1}$ decrease strictly to the Abel value
$g(f)=\lim_{t\uparrow1}\sum_{n\ge0}(-t)^nf(n)$, with
\[
 E_N(f)-H_r^{(b)}2^{-N}<g(f)<E_N(f).
\]
Thus the positivity and Euler remainder extend to independent strides
without requiring $a\ge1$ or $a+b\ge1$.
\end{corollary}
\begin{proof}
The coefficient is
$\iint v^{\lambda n}(1-u^{rn})\dd\sigma_a(v)\dd\tau_b(u)$.
The preceding convexity argument applies with $U=u^r$ and
$V=v^\lambda$, and $\int(1-u^r)\dd\tau_b=H_r^{(b)}$.
For $0<t<1$, geometric summation gives the integrand
$(1+tv^\lambda)^{-1}-(1+tu^rv^\lambda)^{-1}$.
Its absolute value is at most $1-u^r$, so dominated convergence
identifies the Abel limit. The same positive Euler-tail summation
proves the claimed enclosure. This proves an evaluator theorem, not
an assertion about an affine one-crossing finite signed kernel.
\end{proof}

\begin{proposition}[A complex Euler identity and a sharp negative-axis tail]
\label{prop:complexEuler}
For $\Repart z<1/2$, put $q=z/(z-1)$. Then
\begin{equation}\label{eq:complexEuler}
 F_{a,b}(z)=\frac z{1-z}\sum_{k\ge0}d_k^{(1)}q^k,
\end{equation}
and truncation before $k=N$ has error at most
\begin{equation}\label{eq:complexEulerbound}
 \frac{|z|}{|1-z|}\frac{|q|^N}{1-|q|}.
\end{equation}
For $z=-r$, $r>0$, the truncated values increase to $F_{a,b}(-r)>0$,
and the error is strictly less than
\begin{equation}\label{eq:negativeEulerbound}
 r\left(\frac r{1+r}\right)^N\qquad(N\ge1).
\end{equation}
The coefficient one multiplying this last bound is sharp uniformly
over positive $a,b$ at each fixed $r,N$.
\end{proposition}
\begin{proof}
Set $A(t)=\sum_{n\ge0}c_{n+1}t^n=F(t)/t$. The binomial
transform gives
\[
 \sum_{k\ge0}d_k^{(1)}q^k
 =\frac1{1-q}A\left(\frac{-q}{1-q}\right)
\]
initially near zero. The left side is analytic for $|q|<1$ by
$|d_k^{(1)}|<1$. Substitution $q=z/(z-1)$ maps this disk to
$\Repart z<1/2$ and analytic continuation proves the identity.
The geometric tail proves \eqref{eq:complexEulerbound}.
At $z=-r$, $q=r/(1+r)$ and the prefactor is $-q$; since all the
nonzero $d_k$ are negative, monotonicity and the strict bound follow.
For fixed $k$, $d_k\to-1$ as $a\downarrow0$, $b\to\infty$.
Dominated convergence in the geometrically weighted tail yields
exactly $rq^N$, proving sharpness.
\end{proof}

\section{The sharp critical Euler constant}
\label{sec:sharpconstant}
\begin{theorem}[Exact fixed-pair constant on the closed subcritical region]
\label{thm:pairconstant}
For $a,b>0$ with $a+b\le1$, the sequence
$2^N(E_N-g_{a,b})$, $N\ge1$, decreases strictly to zero. In fact
\begin{equation}\label{eq:subcriticaltail}
 2^N(E_N-g_{a,b})
 =\int_0^1\frac{(1-u^2)^N}{1+u^2}\dd\nu_{a,b}(u).
\end{equation}
The least constant $C$ such that
$E_N-g_{a,b}\le C2^{-N}$ for every $N\ge1$ is
\begin{equation}\label{eq:pairconstant}
 C(a,b)=-2g_{a,b}.
\end{equation}
For strict inequalities, the infimum is the same, but is not attained
at a fixed pair because equality occurs at $N=1$.
\end{theorem}
\begin{proof}
For $k\ge1$, the finite compensated moment formula gives
\[
 d_k^{(2)}=-\int_0^1(1-u^2)^k\dd\nu_{a,b}(u).
\]
Its integral is finite: near one, $(1-u^2)^k=O(1-u)$, and
$(1-u)\nu$ is finite. Summing the positive tail proves
\eqref{eq:subcriticaltail}. For $N\ge1$ its integrand is dominated
by $q_1(u)\le2(1-u)$ and decreases strictly to zero for $0<u<1$.
Dominated convergence proves the asserted decrease and limit even
when $\nu$ has infinite mass. The largest scaled error is at $N=1$,
where $E_1=0$, so it equals $-2g_{a,b}$.
\end{proof}

\begin{theorem}[Resolution of the sharp critical constant conjecture]
\label{thm:criticalconstant}
The function $a\mapsto-2g_{a,1-a}$ decreases strictly on $(0,1)$, with
\begin{equation}\label{eq:criticalendpoints}
 \lim_{a\downarrow0}(-2g_{a,1-a})=\frac\pi4+\frac{\log2}{2},
 \qquad
 \lim_{a\uparrow1}(-2g_{a,1-a})=\frac\pi2-1.
\end{equation}
The sharp constant uniform over $0<a<1$, $b=1-a$, and $N\ge1$ is
\begin{equation}\label{eq:Ccritical}
 \boxed{C_{\mathrm{critical}}=\frac\pi4+\frac{\log2}{2}.}
\end{equation}
In particular,
\begin{equation}\label{eq:criticalenclosure}
 E_N-C_{\mathrm{critical}}2^{-N}<g_{a,1-a}<E_N
 \qquad(0<a<1,\ N\ge1).
\end{equation}
No smaller uniform constant works, even if the inequalities are
allowed to be non-strict.
\end{theorem}
\begin{proof}
The strict monotonicity is Theorem~\ref{thm:weightmonotonicity} at
$w=1$, $\rho=1$. Corollary~\ref{cor:weightbounds}, using
$\beta(1)=\pi/4$, $\beta(0)=1/2$, and $\eta_D(1)=\log2$, gives the
endpoint values. The fixed-pair constant is
Theorem~\ref{thm:pairconstant}. Its supremum is therefore exactly
\eqref{eq:Ccritical}. It is not attained by an interior pair, which
proves the strict uniform bound. For any smaller candidate, take
$a$ sufficiently close to zero and $N=1$ to contradict it.
\end{proof}

This proves both parts of the conjecture labeled \texttt{conj:constant}
in \cite{threshold}: strict parameter monotonicity, not merely an
upper bound, and the matching optimal uniform constant. Numerically,
$C_{\mathrm{critical}}\approx1.1319717536774209643$.
Theorem~\ref{thm:pairconstant} and \eqref{eq:universalgaussian} also give
a uniform, not asserted optimal, constant $(1+\sqrt2)/2$ on the whole
closed subcritical triangle.

\section{Strict polylogarithms at arbitrary depth}
\label{sec:depth}
Let $d\ge1$, $\sv=(s_1,\ldots,s_d)$ with $s_j>0$, and $w=\sum s_j$.
Define
\begin{equation}\label{eq:strictdepth}
 F_{\sv}(z)=\sum_{n_1>\cdots>n_d\ge1}
                  \frac{z^{n_1}}{n_1^{s_1}\cdots n_d^{s_d}}
           =\sum_{n\ge d}c_n(\sv)z^n.
\end{equation}
The notation $c_n(\sv)$ in this section generalizes, and does not alter,
the depth-two convention.

\begin{theorem}[Ordered-simplex positive kernel]\label{thm:simplex}
For $z\in\cut$,
\begin{equation}\label{eq:simplex}
 F_{\sv}(z)=z^d\int_{1>x_1>\cdots>x_d>0}
 \frac{W_{\sv}(\xv)}{\prod_{j=1}^d(1-zx_j)}\dd\xv,
\end{equation}
where $x_0=1$ and
\begin{equation}\label{eq:simplexweight}
 W_{\sv}(\xv)=\prod_{j=1}^d
      \frac{\log(x_{j-1}/x_j)^{s_j-1}}{\Gamma(s_j)}.
\end{equation}
The weight is positive, integrable, and has mass
\begin{equation}\label{eq:massdepth}
 M_{\sv}=\prod_{j=1}^d(d-j+1)^{-s_j}=c_d(\sv).
\end{equation}
\end{theorem}
\begin{proof}
Insert $n_j^{-s_j}=\int_0^\infty e^{-n_jt_j}k_{s_j}(t_j)\dd t_j$
and write the positive integer gaps $k_j=n_j-n_{j+1}$, with
$n_{d+1}=0$. Set $T_j=t_1+\cdots+t_j$. Then
$\sum n_jt_j=\sum k_jT_j$ and $n_1=\sum k_j$.
The sums over the $k_j$ are independent geometric series, giving
\[
 \int_{\R_+^d}\prod_{j=1}^d k_{s_j}(t_j)
          \prod_{j=1}^d\frac{ze^{-T_j}}{1-ze^{-T_j}}\dd\tv.
\]
All exchanges are justified by absolute convergence for $|z|<1$.
Now set $x_j=e^{-T_j}$. The Jacobian is $\prod x_j$, exactly
canceling the product in the numerators. This yields
\eqref{eq:simplex}. For its mass, $\prod x_j=e^{-\sum_{i=1}^d(d-i+1)t_i}$,
so independent gamma integrals give \eqref{eq:massdepth}.
Finite mass and denominator separation on compact slit sets prove
holomorphic continuation as in Lemma~\ref{lem:double}.
\end{proof}

\subsection{A positive generalized Stieltjes measure}
Let $\Delta_{d-1}=\{p_j\ge0:\sum p_j=1\}$, with uniform Dirichlet
probability measure (density $\Gamma(d)$ in its usual $d-1$ coordinates).
For $d=1$ this is the point mass at $p_1=1$.

\begin{theorem}[Depth-$d$ generalized Stieltjes representation]
\label{thm:Stieltjesrepresentation}
There is a finite positive measure $\mu_{\sv}$ on $(0,1)$, with mass
$M_{\sv}$ and positive mass on every nonempty interior interval, such that
\begin{equation}\label{eq:Stieltjes}
 \frac{F_{\sv}(z)}{z^d}
 =\int_0^1(1-zu)^{-d}\dd\mu_{\sv}(u)\qquad(z\in\cut).
\end{equation}
Explicitly, $\mu_{\sv}$ is the pushforward of
$W_{\sv}(\xv)\dd\xv$ times the uniform Dirichlet probability law under
$u=\sum_{j=1}^dp_jx_j$.
\end{theorem}
\begin{proof}
For $\Repart A_j>0$, independent Laplace integrals followed by
$r_j=tp_j$ give the classical parameter identity
\begin{equation}\label{eq:Feynman}
 \prod_{j=1}^dA_j^{-1}
 =\Gamma(d)\int_{\Delta_{d-1}}
              \left(\sum_{j=1}^dp_jA_j\right)^{-d}\dd\pv.
\end{equation}
Apply it with $A_j=1-zx_j$ first for $\Repart z<1$. Absolute
integrability against the finite positive weight justifies the
iterated integral. The proposed pushforward then proves
\eqref{eq:Stieltjes}; analytic continuation gives the whole slit plane.
Since every $x_j$ lies in $(0,1)$, so does $u$.
For any open interval $I\subset(0,1)$, choose all strictly ordered
$x_j$ in $I$. That set has positive $W$-mass and its every convex
combination lies in $I$. Hence $\mu_{\sv}(I)>0$.
\end{proof}

The notion of exact generalized Stieltjes order is standard; see
Karp and Prilepkina \cite{karp-prilepkina}. We prove the needed minimality
from a decay obstruction rather than invoking an exact-order criterion.
The recent star-polylogarithm representation of Li \cite{li-star}
concerns weakly decreasing indices and different normalizations. It is
not a contradiction to the exact order found here for the strictly
nested function divided by $z^d$.

\subsection{A negative-axis asymptotic that forces exact order}
\begin{theorem}[Universal logarithmic leading term]\label{thm:negativeasymptotic}
For every positive index vector $\sv$,
\begin{equation}\label{eq:negativeasymptotic}
 (-1)^dF_{\sv}(-x)\sim\frac{(\log x)^w}{\Gamma(w+1)}
 \qquad(x\to+\infty).
\end{equation}
The leading coefficient depends only on the total weight, not on the
individual positive orders.
\end{theorem}
\begin{proof}
Put $x=e^L$ and $\ell(y)=(1+e^{-y})^{-1}$. The Mellin representation
in the proof of Theorem~\ref{thm:simplex} gives the exact positive formula
\begin{equation}\label{eq:logistic}
 (-1)^dF_{\sv}(-e^L)
 =\frac1{\prod_j\Gamma(s_j)}\int_{\R_+^d}
       \prod_j t_j^{s_j-1}\prod_{j=1}^d\ell(L-T_j)\dd\tv.
\end{equation}
Rescale $t_j=Lv_j$ for $L\ge1$. Except on a set of measure zero, the
product of logistic factors tends to the indicator of $\sum v_j<1$,
because the cumulative sums increase with $j$.
On $\sum v_j\le1$ the product is bounded by one. On $r=\sum v_j>1$,
the last factor is at most $e^{-L(r-1)}\le e^{-(r-1)}$, while all
others are at most one. These bounds, multiplied by
$\prod v_j^{s_j-1}$, are integrable: radial integration reduces the
outside bound to a constant times $\int_1^\infty r^{w-1}e^{-(r-1)}\dd r$.
Dominated convergence now gives
\[
 \lim_{L\to\infty}\frac{(-1)^dF_{\sv}(-e^L)}{L^w}
 =\frac1{\prod_j\Gamma(s_j)}
   \int_{v_j>0,\,\sum v_j<1}\prod_jv_j^{s_j-1}\dd\boldsymbol v
 =\frac1{\Gamma(w+1)}.
\]
\end{proof}

\begin{definition}
A generalized Stieltjes function of order $\alpha>0$ is a function
on $x>0$ representable as
\[
 f(x)=c+\int_{[0,\infty)}(x+t)^{-\alpha}\dd\lambda(t),
 \qquad c\ge0,\quad\lambda\ge0,
\]
with $\int(1+t)^{-\alpha}\dd\lambda(t)<\infty$.
Its exact order is the infimum of the possible positive orders.
\end{definition}

\begin{theorem}[Exact generalized Stieltjes order equals strict depth]
\label{thm:exactorder}
The function
\begin{equation}\label{eq:Hdepth}
 H_{\sv}(x)=\frac{F_{\sv}(-x)}{(-x)^d}\qquad(x>0)
\end{equation}
has exact generalized Stieltjes order $d$. It has no positive
representation of any order $0<\alpha<d$, even when the representing
measure is allowed on the entire half-line rather than a compact interval.
\end{theorem}
\begin{proof}
Formula~\eqref{eq:Stieltjes} represents $H_{\sv}$ as
$\int(1+xu)^{-d}\dd\mu_{\sv}(u)$. With $t=1/u$ and the weighted
pushforward $\dd\lambda(t)=u^{-d}\dd\mu_{\sv}(u)$, this is a generalized
Stieltjes representation of order $d$. The required integrability follows
from $\int(1+u)^{-d}\dd\mu_{\sv}<\infty$.

By Theorem~\ref{thm:negativeasymptotic},
\begin{equation}\label{eq:Hdecay}
 H_{\sv}(x)\sim\frac{(\log x)^w}{\Gamma(w+1)x^d}.
\end{equation}
Suppose it had an order-$\alpha$ representation for $\alpha<d$.
Since $H_{\sv}(x)\to0$, the constant term is zero. The nonzero positive
measure must assign a finite positive mass $m$ to some $[0,R]$;
local finiteness follows from the defining weighted integrability.
Then $H_{\sv}(x)\ge m(x+R)^{-\alpha}$. But \eqref{eq:Hdecay} implies
$x^\alpha H_{\sv}(x)\to0$, a contradiction.
\end{proof}

\section{The exact coefficient-normalization threshold}
\label{sec:normalization}
For $\alpha>0$, let $(\alpha)_n=\Gamma(\alpha+n)/\Gamma(\alpha)$ and
set
\begin{equation}\label{eq:normalization}
 M_n^{(\alpha)}=\frac{n!}{(\alpha)_n}c_{n+d}(\sv),\qquad n\ge0.
\end{equation}
\begin{theorem}[Positive Hausdorff normalization if and only if $\alpha\ge d$]
\label{thm:normalization}
The sequence $(M_n^{(\alpha)})_{n\ge0}$ is a positive Hausdorff moment
sequence on $[0,1]$ exactly when $\alpha\ge d$.
At $\alpha=d$, its measure is $\mu_{\sv}$. At $\alpha>d$, a measure
is obtained by replacing $U$ by $BU$, where $U$ has normalized law
$\mu_{\sv}/M_{\sv}$ and $B\sim\operatorname{Beta}(d,\alpha-d)$ is
independent, then multiplying the law by $M_{\sv}$.
For every $\alpha\ge d$, all finite alternating differences are
strictly positive and all finite principal Hankel matrices are
positive definite.
\end{theorem}
\begin{proof}
Expand \eqref{eq:Stieltjes} inside the disk:
\[
 c_{n+d}(\sv)=\frac{(d)_n}{n!}\int_0^1u^n\dd\mu_{\sv}(u).
\]
This proves the assertion at $\alpha=d$. For $\alpha>d$,
$\E[B^n]=(d)_n/(\alpha)_n$, giving the product-measure construction.
The measure assigns positive mass to every interior interval. Therefore
\[
 \sum_{j=0}^k(-1)^j\binom kj M_{n+j}^{(\alpha)}
 =\int u^n(1-u)^k\dd\mu^{(\alpha)}(u)>0,
\]
and the quadratic-form argument for a nonzero polynomial gives strict
positive definiteness of every finite Hankel matrix.

Conversely, suppose $\alpha<d$ and these numbers were moments of a
positive finite measure $\lambda$ on $[0,1]$. Summing its moments would give
\[
 \frac{F_{\sv}(z)}{z^d}=\int_0^1(1-zu)^{-\alpha}\dd\lambda(u)
\]
in the disk, and hence on the slit plane. A possible atom at zero is
a nonnegative constant; the remainder, after $t=1/u$, is a generalized
Stieltjes function of order $\alpha$. This contradicts
Theorem~\ref{thm:exactorder}.
\end{proof}

\begin{example}[A finite witness against the unnormalized order-one claim]
For $\sv=(1,1)$, depth $d=2$, and $\alpha=1$,
\[
 M_0^{(1)}=\frac12,\qquad M_1^{(1)}=\frac12,
 \qquad M_2^{(1)}=\frac{11}{24}.
\]
Their second alternating difference is
$M_0^{(1)}-2M_1^{(1)}+M_2^{(1)}=-1/24<0$.
In contrast, the correctly normalized coefficients
$c_{n+2}/(n+1)$ are positive Hausdorff moments at every positive pair
of real indices. An omitted normalization is thus a checkable error,
not a harmless convention.
\end{example}

\begin{corollary}[Integral and differential normalization identities]
\label{cor:normalizingidentities}
Put $A_{\sv}(z)=F_{\sv}(z)/z^d$, with its removable value at zero, and
\[
 \mathcal M_{\sv}(z)=\sum_{n\ge0}
       \frac{c_{n+d}(\sv)}{\binom{n+d-1}{d-1}}z^n
       =\int_0^1\frac{\dd\mu_{\sv}(u)}{1-zu}.
\]
For $d\ge2$ and $z\in\cut$,
\begin{align}
 \mathcal M_{\sv}(z)
 &=(d-1)\int_0^1(1-t)^{d-2}A_{\sv}(tz)\dd t,
 \label{eq:normalizingintegral}\\
 A_{\sv}(z)
 &=\frac1{(d-1)!}\frac{\dd^{d-1}}{\dd z^{d-1}}
                       \bigl[z^{d-1}\mathcal M_{\sv}(z)\bigr].
 \label{eq:normalizingderivative}
\end{align}
\end{corollary}
\begin{proof}
The beta integral gives
$(d-1)\int_0^1t^n(1-t)^{d-2}\dd t=1/\binom{n+d-1}{d-1}$.
This proves \eqref{eq:normalizingintegral} by coefficient comparison.
Differentiation multiplies the $n$th coefficient by the reciprocal
factor, proving \eqref{eq:normalizingderivative}. Both sides are
holomorphic on $\cut$: for a compact slit set, its products with
$t\in[0,1]$ form a compact subset of the same slit plane. Analytic
continuation therefore establishes the stated domain.
\end{proof}

\begin{remark}[Analytic order is not arithmetic depth]
The exact order $d$ just proved is a positivity and decay invariant of
this analytic function with its specified normalization. It does not
prove that a special value cannot be reduced to lower-depth periods,
that numerical constants are independent, or that a motivic depth
filtration has a particular rank. These are distinct questions.
\end{remark}

\section{Reproducible exact certificates and independent tests}
\label{sec:certificates}
\subsection{Fractional powers without floating-point assumptions}
The standard-library script \texttt{code/certify.py} uses integer and
rational arithmetic only. For $p/q>0$ and integer $n\ge1$, put $B=10^{80}$
and compute
\[
 r=\left\lfloor(n^pB^q)^{1/q}\right\rfloor.
\]
Before using the result, it verifies both integer inequalities
\begin{equation}\label{eq:rootcertificate}
 r^q\le n^pB^q<(r+1)^q.
\end{equation}
Except in the exact-root case, these imply
$B/(r+1)<n^{-p/q}<B/r$. The exact-root case is handled separately.
Endpoints are rounded outward onto $B^{-1}\mathbb Z$ using integer
division. Positive harmonic sums and products then enclose $c_n$.
No floating-point root or transcendental sign is accepted as a proof.

The finite Euler sum is formed using the exact weights
\begin{equation}\label{eq:finiteweights}
 E_N=2^{-N}\sum_{n=0}^{N-1}(-1)^n c_{2n+1}
                         \sum_{j=n+1}^{N}\binom Nj.
\end{equation}
Substitution of \eqref{eq:dk} proves this identity by the finite relation
$\sum_{k=n}^{N-1}2^{N-k-1}\binom kn=\sum_{j=n+1}^N\binom Nj$,
which follows from Pascal's recurrence. We use $N=192$ and the proved
rational tail constant $2$ from Theorem~\ref{thm:Eulerall}. The finite
centers at fractional orders are algebraic numbers enclosed by rational
intervals; they are not mislabeled as rational exact values.

\begin{center}
\small
\begin{tabular}{@{}c c l@{}}
\toprule
$a$ & $b$ & $g_{a,b}$ (displayed approximation only)\\
\midrule
$1/10$ & $1/10$ & $-0.490480773154378134692711376810$\\
$1/4$ & $1/4$ & $-0.464611276461282330148775822252$\\
$2/5$ & $1/5$ & $-0.424200725534215249048994001121$\\
$1/3$ & $1/3$ & $-0.446085363511268827586991313068$\\
$1/10$ & $9/10$ & $-0.529761932285771263094629365374$\\
$1/2$ & $1/2$ & $-0.403936987717013579833068109579$\\
$9/10$ & $1/10$ & $-0.306037812999015482508621228917$\\
$1$ & $1$ & $-0.272198261287950266312586112280$\\
$2$ & $3$ & $-0.093247950371174729821546998597$\\
\bottomrule
\end{tabular}
\end{center}
The exact JSON records both rational endpoints and outward decimal
summaries. Every interval has width below $3.187\times10^{-58}$.
For a shorter readable certified example,
\begin{align*}
 -0.464611276461282330148775822252
 &<g_{1/4,1/4}\\
 &<-0.464611276461282330148775822251.
\end{align*}
This lies strictly below the old finite-measure threshold $a+b=1$.
It is a value enclosure, not a conjectured closed form.

\subsection{Angular brackets with rational Chebyshev evaluation}
For rational $0<\rho<1$ and $-1<c<1$, the exact identity is
\begin{equation}\label{eq:chebyshev}
 \frac{\Impart F_{a,b}(\rho e^{\ii\theta})}{\sin\theta}
 =\sum_{n\ge1}c_n\rho^nU_{n-1}(c),\qquad c=\cos\theta.
\end{equation}
Use $c_n\le n$ and $|U_{n-1}(c)|\le n$. For $j=N+1$, an explicit
absolute tail bound is
\begin{equation}\label{eq:angulartail}
 \rho^j\left\{\frac{j^2}{1-\rho}
       +\frac{2j\rho}{(1-\rho)^2}
       +\frac{\rho(1+\rho)}{(1-\rho)^3}\right\}.
\end{equation}
This is the exact sum of $\sum_{n\ge j}n^2\rho^n$.

For $c=P/Q$, the program computes integer numerators
$U_k(P/Q)=V_k/Q^k$, where
$V_0=1$, $V_1=2P$, and $V_k=2PV_{k-1}-Q^2V_{k-2}$.
Each rational summand is rounded outward separately. This avoids
unstable interval recursion in the Chebyshev polynomial itself.
The verifier uses $N=400$, checks a strictly negative enclosure at
the lower cosine endpoint and a strictly positive one at the upper
endpoint, and checks that both lie between $\rho/2$ and $\rho$.
The analytic theorem, not the finite root search, supplies uniqueness.

\begin{center}
\small
\begin{tabular}{@{}c c r r r@{}}
\toprule
$a$ & $b$ & $\eta(1/4)$ & $\eta(1/2)$ & $\eta(3/4)$\\
\midrule
$1/10$ & $1/10$ & $0.9290066544$ & $0.9320671593$ & $0.9391522820$\\
$1/4$ & $1/4$ & $0.8330769037$ & $0.8376577074$ & $0.8478255853$\\
$2/5$ & $1/5$ & $0.7966236703$ & $0.8012312684$ & $0.8112578998$\\
$1/3$ & $1/3$ & $0.7848252542$ & $0.7893512974$ & $0.7991218117$\\
\bottomrule
\end{tabular}
\end{center}
These are numerical summaries. The twelve accepted rational cosine
brackets each have width $5\times10^{-12}$. No trigonometric rounding
is involved in their acceptance. No unit-radius bracket is claimed;
unit radius is covered by the endpoint proof in Lemma~\ref{lem:root}.

\subsection{Recorded evidence and replay}
The exact verifier checks 3,915 integer-root inequalities, all nine
Euler enclosures, and all twelve angular brackets. At integer orders
it also verifies 728 strictly positive finite differences and 96
strictly positive Hankel determinants for eight index vectors through
depth four, together with the exact counterexample $-1/24$ above.
A further 288 exact coefficient checks cover independent affine strides
at integer orders, including the noninteger denominator step $1/2$.
The separate SymPy script checks seven rational identities used in
the proofs, including the root derivative and the beta-kernel derivative.
These finite instances are tests, not substitutes for the all-depth
arguments.

Independent generalized Gauss--Laguerre quadrature uses
$2^{-a}$ times a product of gamma probability laws after
$x=e^{-S/2}$, $y=xe^{-T}$. Its 128-node comparison against the Euler
centers has largest observed absolute discrepancy
$3.109\times10^{-15}$ in the recorded run. This is a floating-point
diagnostic, not a certified quadrature error. Six fixed-weight grids
with nine points each show strict decrease, again only as diagnostics.

From the package root, the exact replay is
\begin{verbatim}
python code/certify.py --check
\end{verbatim}
It requires only the Python standard library. Optional diagnostics and
symbolic tests are reproduced by
\begin{verbatim}
python code/diagnostics.py
python code/symbolic_checks.py
\end{verbatim}
The first command regenerates rational bracket proposals; running
\texttt{certify.py} afterward independently accepts or rejects them.
For the frozen replay, do not regenerate proposals first.
The file manifest checks byte integrity, not mathematical truth.

\section{Audit conclusions and integration guidance}
\label{sec:audit}
\subsection{What is corrected, and what is only strengthened}
No contradiction to the \emph{stated} finite signed threshold or the
previous critical endpoint theorem was found. Those results remain
valid. The subcritical construction is a different representation, not
a repair of a false finite-measure theorem. The precise status changes
are that the prior angular-uniqueness and sharp-critical-constant
conjectures now have proofs, and the separate subcritical slit-plane
question has a positive answer.

There are four concrete hazards to correct in any integration or
extrapolation of earlier material.

\paragraph{Infinite mass must remain compensated.}
Below the threshold, splitting the new positive integral into separate
masses is invalid. Use \eqref{eq:compensatedmoments} or the finite measure
$\omega$, never a fictitious finite signed measure. The threshold
$a+b\ge1$ and the first-difference threshold $a+b\le1$ are complementary,
not competing claims about the same measure.

\paragraph{Boundary values are not ordinarily convergent sums.}
For $w<1$ the coefficients grow, and on $w=1$ they approach a nonzero
constant. The new Euler theorem evaluates analytic/Abel values in
these regimes. It must not be described as acceleration of an
ordinarily convergent boundary series there.

\paragraph{The singular angular endpoint needs its own proof.}
At $\rho=1$ below threshold, the limiting integral
$\int(1-u)^{-1}\dd\omega$ is infinite. An argument using dominated
convergence against that expression would be invalid. The negative-tail
bound by $2\omega((m,1))$ in Lemma~\ref{lem:root} is the required
replacement and also proves existence, not just uniqueness conditional
on a zero being present.

\paragraph{Stieltjes normalizations and depth conventions matter.}
For strict depth $d$, the order-one moment sequence is
$c_{n+d}/\binom{n+d-1}{d-1}$, not generally $c_{n+d}$. The exact
counterexample in Section~\ref{sec:normalization} shows that dropping
the denominator can falsify even a three-term positivity test. A star
polylogarithm theorem must not be transferred to strict nesting without
checking this change. Exact analytic Stieltjes order is not an
arithmetic non-reducibility theorem for its special values.

\subsection{Suggested insertion points}
The complete report can be placed at
\nolinkurl{Analysis/Polylogarithms/docs/reports/compensated-polylogarithms/}.
The supplied insertion fragment uses the prefix \texttt{compensated:}.
It summarizes the new theorems and points back to the full proofs.
The current main file inputs \texttt{chapters/05-certified-computation.tex};
that chapter and the signed-kernel discussion are natural insertion
points for the compensated and Euler results. The zero-geometry chapter
and final research-programme chapter should cross-reference the two
resolved conjectures. The all-depth Stieltjes theorem fits with the
foundational iterated-integral material.

Integration should preserve the historical record: label the original
conjectures as resolved by this continuation, do not rewrite them as
if they had already been proved in the earlier report. Keep the
supercritical normalized-radius conjecture, $S_6$, and numerical period
independence questions explicitly open. A stale or differently organized
chapter path in an earlier integration note should be mapped to the
current main-file inputs rather than copied blindly.

\section{Further research: precise next questions}
\label{sec:future}
\subsection{The remaining normalized-radius phase diagram}
We have proved strict increase for $a+b\le1$. The earlier conjecture
for integer $a\ge2$ and finite $b>0$, and the sign transition at $a=1$,
are not settled by this argument. A first objective is a sharp description
of the set of $(a,b)$ for which the radial sign changes as $\rho$ varies.
The small-radius coefficient in \eqref{eq:smallradius} is an exact local
test, but a sign for that coefficient does not imply a global sign.
The positive variance mechanism here works because the factor is
$1-z$, equivalently because $d(u)-d(m)=(u-m)(u-1)$; an interior
crossing factor $1-rz$ need not have the same sign structure.

\subsection{Optimal Euler constants beyond the critical line}
The exact fixed-pair constant is now known on $a+b\le1$. A sharper
uniform constant on the whole subcritical triangle is reduced to
understanding the envelope
\[
 C_*(w)=\beta(w)+2^{-w}\eta_D(w),\qquad0<w\le1,
\]
since the fixed-weight supremum occurs as $a\downarrow0$.
The bound $(1+\sqrt2)/2$ is explicit but is not asserted optimal.
On the full positive quadrant, the all-order Euler bound $H_2^{(b)}$
is valid, but scaled-error monotonicity in $N$ is not proved here.
Finding the actual universal Gaussian Euler constant remains a separate
optimization problem. Uniform error asymptotics near $w=1$ would also
clarify how infinite compensation crosses over to a finite critical
measure.

\subsection{Higher-depth angular geometry}
The exact Stieltjes order suggests a natural angular-count question,
but order-$d$ positivity alone does not prove it. There is a rigorous
small-radius starting point.
\begin{proposition}[The small-radius angular count]\label{prop:smallcount}
For a fixed positive index vector of depth $d$, and every sufficiently
small $\rho>0$, $\Impart F_{\sv}(\rho e^{\ii\theta})$ has exactly
$d-1$ simple zeros in $0<\theta<\pi$.
\end{proposition}
\begin{proof}
After division by $\rho^d\sin\theta$, its analytic coefficient expansion
is
\[
 M_{\sv}U_{d-1}(c)+\sum_{n>d}c_n(\sv)\rho^{n-d}U_{n-1}(c),
 \qquad c=\cos\theta.
\]
The remainder tends to zero in $C^1([-1,1])$ as $\rho\downarrow0$.
For instance $c_n\le n^{d-1}$, $|U_{n-1}|\le n$, and
$|U_{n-1}'|\le n^3$ on this interval. The latter bound follows by
writing $U_{n-1}(\cos\theta)$ as a finite sum of cosines and using
$|\sin(k\theta)|\le k|\sin\theta|$. These polynomial bounds give
uniform convergence with the first derivative. The polynomial
$U_{d-1}$ has precisely $d-1$ simple roots $\cos(j\pi/d)$ in $(-1,1)$
and no endpoint roots. Choose small disjoint neighborhoods of these
roots where its derivative has a fixed nonzero sign. A sufficiently
small $C^1$ perturbation has one simple root in each neighborhood
and none on their compact complement. The case $d=1$ has no roots.
\end{proof}

\begin{conjecture}[Full-radius strict-depth angular count]
\label{conj:depthangular}
For every positive index vector of depth $d$ and every $0<\rho\le1$,
there are exactly $d-1$ simple zeros of
$\Impart F_{\sv}(\rho e^{\ii\theta})$ in $0<\theta<\pi$.
\end{conjecture}
At depth one, the kernel $z/(1-zu)$ has strictly positive imaginary
part in the upper half-plane, so there are no angular zeros.
The depth-two assertion is proved in this article. The arbitrary-depth
assertion is a proposed research conjecture, supported here only by
the small-radius proposition and the depth-one and depth-two cases;
no exhaustive numerical higher-depth evidence is claimed.

\subsection{Constructive measures and sharper asymptotics}
The pushforward measure $\mu_{\sv}$ has a concrete positive construction,
but an explicit one-variable density would improve quadrature and
Pad\'e computations. Conditional Dirichlet averages should permit a
piecewise-polynomial kernel in the $x_j$ variables before integration
against the logarithmic weight. It would be useful to derive endpoint
asymptotics of this density with uniform constants as some $s_j$ approach
zero, and to sharpen \eqref{eq:negativeasymptotic} beyond the leading term.
The boundary layer $\sum t_j\approx\log x$ interacts with the lower
faces $t_j\approx0$, so an all-orders expansion must track individual
indices rather than only the total weight.

\subsection{Colored periods, exact identities, and formalization}
Other outer colors on the unit circle can destroy the positive sign
of the double kernel. The right question is to characterize parameter
and color sectors where a positive real kernel survives, rather than
reuse a Gaussian sign argument formally. The normalization identities
\eqref{eq:normalizingintegral}--\eqref{eq:normalizingderivative} may
supply stable new coordinates for evaluating such colored functions.
They do not resolve the existing $S_6$ period identity, whose finite
mixed-reduction certificate remains a concrete arithmetic target.

The finite formalization core is small: the integer-root interval
operations, binomial Euler weights, the rational root-kernel derivatives,
and the Chebyshev sign checker. A subsequent analytic formalization can
separate the compensated measure, the beta stochastic-order lemma,
compact-slit continuation, and the large-negative-axis dominated
convergence argument. The current scripts are reproducible exact
arithmetic, not a Lean development.

\section*{Conclusion}
\addcontentsline{toc}{section}{Conclusion}
The finite signed-measure boundary is real, but it is not the boundary
of positive structure. Below it, compensation produces a finite positive
first-difference measure and restores the analytic geometry that the old
representation could not reach. This proves subcritical angular
uniqueness, slit-plane nonvanishing, and normalized-radius increase.
A different positivity mechanism---beta likelihood-ratio order at fixed
total weight---settles the sharp critical Euler constant. Positive
transport extends the evaluator to every positive pair, while ordered
simplex kernels identify an exact Stieltjes order at all strict depths.
Each statement has its own explicit scope, and the remaining period and
supercritical phase questions retain their unresolved status.

\clearpage
\addcontentsline{toc}{section}{References}
\begin{thebibliography}{9}
\bibitem{repo}
ProveIt Contributors,
\emph{Polylogarithms and their Arithmetic Bridges}, canonical manuscript,
repository \texttt{VladimirReshetnikov/ProveIt}, revision \basecommit.
\url{https://github.com/VladimirReshetnikov/ProveIt/tree/a0a90ef31877f98be437191c48b46f02d5456867/Analysis/Polylogarithms/docs/manuscript}.

\bibitem{threshold}
ProveIt Contributors,
\emph{The Sharp Real-Order Threshold for Harmonic Polylogarithms:
Endpoint atoms, zero geometry, and certified fractional-order identities},
9 October 2026. Incoming archive
\texttt{ProveIt\_RealOrder\_Threshold.zip}; source
\texttt{real\_order\_threshold.tex}.
Conjectures \texttt{conj:subcritical} and \texttt{conj:constant} are the
specific conjectures resolved here. The accessible Library source was
read and hashed; see \texttt{PROVENANCE.json}.
\url{https://github.com/VladimirReshetnikov/ProveIt/blob/a0a90ef31877f98be437191c48b46f02d5456867/docs/incoming/ProveIt_RealOrder_Threshold.zip}.

\bibitem{dlmf-hurwitz}
NIST Digital Library of Mathematical Functions,
\S25.11, \emph{Hurwitz Zeta Function}, especially the integral
representations and the regularized formula 25.11.27.
Accessed during this continuation, October 2026.
\url{https://dlmf.nist.gov/25.11}.

\bibitem{dlmf-polylog}
NIST Digital Library of Mathematical Functions,
\S25.12, \emph{Polylogarithms}. Standard principal-branch conventions and
integral representations. Accessed October 2026.
\url{https://dlmf.nist.gov/25.12}.

\bibitem{karp-prilepkina}
D. Karp and E. Prilepkina,
\emph{Generalized Stieltjes transforms: basic aspects},
arXiv:1111.4271 (2011). Published as
\emph{Generalized Stieltjes functions and their exact order},
Journal of Classical Analysis \textbf{1} (2012), 53--74.
\url{https://arxiv.org/abs/1111.4271}.

\bibitem{li-star}
J. Li,
\emph{The derived set of multiple star polylogarithms},
arXiv:2609.03653v1, 3 September 2026. Cited for a nearby but distinct
star-polylogarithm Hausdorff--Stieltjes construction, not as a premise of
our strict-depth proofs.
\url{https://arxiv.org/abs/2609.03653}.
\end{thebibliography}
\end{document}
